\subsection{The coin-mixture prior}\label{sec:coins}

The simplest exchangeable scaling family is a mixture of
independent-answer models. It gives a complete calculus exercise and
a clean example of the boundary layer at $t = 0$. Throughout this
section the answer width $m$ and the component family are fixed while
$n \to \infty$.

\subsubsection{Categorical formulation}

Let $Z$ take values in a finite set of components indexed by $r$, with
$\Pr(Z = r) = w_r$ summing to one. The mixture defines a joint
distribution for the component $Z$ and the random map $\psi$, the
component being sampled once for the entire map rather than once per
question. Conditional on $Z = r$ the answers at distinct question
indices are independent with categorical distribution $\nu_r$ on the
$A$ possible answers:
%
\begin{equation}
  \nu_r : \mathcal A \to [0,1],
  \qquad
  \sum_{a\in\mathcal A}\nu_r(a) = 1,
  \qquad
  \Pr\big(\psi(q_i) = a \mid Z = r\big) = \nu_r(a).
  \label{eq:coincomponent}
\end{equation}
%
Thus $\nu_r$ acts on \emph{one possible answer} $a$, not on a whole
map: it returns the probability that component $r$ assigns to that
answer at any question. For $m=1$ a Bernoulli component with bias
$\theta_r$ has $\nu_r(1) = \theta_r$; for $m = 2$ a component could
instead be the table $\nu_r(00),\ldots,\nu_r(11)
= \tfrac12,\tfrac14,\tfrac18,\tfrac18$, a distribution on the four
two-bit answers that need not factor into independent bits. The joint
law is
%
\begin{equation}
  \Pr(Z = r, \psi = \phi_j)
  = w_r\prod_{i=1}^{Q}\nu_r\big(\phi_j(q_i)\big),
  \qquad
  p_j = \sum_r w_r\prod_{i=1}^{Q}\nu_r\big(\phi_j(q_i)\big),
  \label{eq:coinprior}
\end{equation}
%
the product being the map's probability conditional on component $r$
and the outer sum averaging those likelihoods. Write
$\nu_{\rm mix} := \sum_r w_r\nu_r$ for the one-question mixture, so
that
%
\begin{equation}
  h_0 = H(\nu_{\rm mix}),
  \qquad
  h_{\rm cond} := \sum_r w_r H(\nu_r)
  \label{eq:coinentropies}
\end{equation}
%
are the entropy of one answer before anything is known and the
entropy per question once the component is known.

Because the construction treats all questions alike, $H(\psi(\mathcal S))$ takes the same value $G_k$ for every $\mathcal S$ of size $k$.
Applying the chain rule to the joint distribution of $Z$ and
$\psi(\mathcal S)$ in both orders,
%
\begin{equation}
  H\big(Z,\psi(\mathcal S)\big)
  = H(Z) + H\big(\psi(\mathcal S)\mid Z\big)
  = H\big(\psi(\mathcal S)\big) + H\big(Z\mid\psi(\mathcal S)\big),
  \label{eq:coinchain}
\end{equation}
%
and using conditional independence,
$H(\psi(\mathcal S)\mid Z) = k\,h_{\rm cond}$, the two expansions
rearrange into the exact decomposition
%
\begin{equation}
  G_k = k\,h_{\rm cond} + I_k,
  \qquad
  I_k := I\big(Z;\psi(\mathcal S)\big),
  \qquad
  \gamma_k = h_{\rm cond} + (I_{k+1}-I_k),
  \label{eq:coindecomp}
\end{equation}
%
in which $I_k$ is the mutual information between the selected
component and the first $k$ answers. The correction has the
conditional-information form
$I_{k+1}-I_k = I(Z;\psi(q)\mid\psi(\mathcal S)) \ge 0$, exactly the
additional information that the next answer reveals about which
component was selected. Substituting into \eqref{eq:finitelev} gives
the exact finite-$n$ leverage
%
\begin{equation}
  \langle L_{n,\ell}\rangle
  = \frac{Q\,H(\nu_{\rm mix})
    - (Q-\ell)\big[h_{\rm cond} + I_{\ell+1}-I_\ell\big]}
    {\ell\,h_{\rm cond} + I_\ell},
  \label{eq:coinfinitelev}
\end{equation}
%
and, provided $h_{\rm cond} > 0$, the inequalities $I_\ell \ge 0$ and
$I_{\ell+1}-I_\ell \ge 0$ show that the finite curve lies below its
limiting hyperbola at the same $t = \ell/Q$.

For finitely many components $0 \le I_k \le H(Z) < \infty$: the
information required to identify the component is subextensive.
Discard any zero-weight components; if the remaining distinct
components are identifiable and mutually absolutely continuous, then
for data generated from a positive-weight component $r$ and any
positive-weight competitor $s$, almost surely under
$\nu_r^{\otimes\infty}$,
%
\begin{equation}
  \frac1k\log_2
  \frac{\Pr(Z = s\mid\psi(\mathcal S_k))}
       {\Pr(Z = r\mid\psi(\mathcal S_k))}
  \qquad
  (\mathcal S_1 \subset \mathcal S_2 \subset \cdots,
   \;|\mathcal S_k| = k)
  \;\longrightarrow\;
  -D_{\rm KL}(\nu_r\Vert\nu_s) \;<\; 0,
  \label{eq:coinodds}
\end{equation}
%
with $D_{\rm KL}$ in base-two logarithms. With a support mismatch the
divergence may be $+\infty$, and one revealing symbol can instead set
the wrong component's posterior probability to zero. Either way the
posterior predictive converges to the true $\nu_r$, and bounded
convergence gives $\gamma(t) = h_{\rm cond}$ for every fixed
$0 < t < 1$. Then $g(t) = t\,h_{\rm cond}$ and \eqref{eq:master}
collapses to
%
\begin{equation}
  L(t) = 1 + \frac ct,
  \qquad
  c = \frac{H(\nu_{\rm mix}) - h_{\rm cond}}{h_{\rm cond}} \;\ge\; 0,
  \label{eq:coinhyperbola}
\end{equation}
%
the inequality being the concavity of entropy,
$H(\sum_r w_r\nu_r) \ge \sum_r w_r H(\nu_r)$. Equality holds exactly
when all positive-weight components share the same categorical
distribution once duplicates are merged, the effective one-coin model
with $L \equiv 1$. If instead $h_{\rm cond} = 0$, as for a nontrivial
mixture of deterministic coins, the map entropy is $O(1)$ and bounded
by $H(Z)$ rather than extensive in $Q$; the normalization of
Section~\ref{sec:tdlcalc} degenerates and
\eqref{eq:coinhyperbola} is not a finite thermodynamic curve.

The similarity to the fixed-$p$ spike is now explicit. In both cases
a global latent variable is inferred within $o(Q)$ questions, leaving
a constant positive entropy density afterwards. The macroscopic
profile sees only the jump from $h_0$ down to $h_{\rm cond}$, and
equation~\eqref{eq:boundarylayer} of
Appendix~\ref{app:endpoints} turns that jump into $1/t$.

\subsubsection{The binary coin and exact finite formulas}

For $m = 1$ write $\nu_r = \operatorname{Bernoulli}(\theta_r)$ and
denote the truth-table Hamming weight by
$|\phi_j| := \sum_i\phi_j(q_i)$, so that
$p_j = \sum_r w_r\theta_r^{|\phi_j|}(1-\theta_r)^{Q-|\phi_j|}$. After
$\ell$ observations containing $k$ ones, let
%
\begin{equation}
  P_{\ell,k} := \sum_r w_r\,\theta_r^k(1-\theta_r)^{\ell-k}.
  \label{eq:coinPlk}
\end{equation}
%
Every particular binary string of length $\ell$ and weight $k$ has
this probability, so
%
\begin{equation}
  G_\ell = -\sum_{k=0}^{\ell}\binom{\ell}{k}
  P_{\ell,k}\log_2 P_{\ell,k},
  \qquad
  w_r(\ell,k)
  = \frac{w_r\theta_r^k(1-\theta_r)^{\ell-k}}{P_{\ell,k}},
  \label{eq:coinGl}
\end{equation}
%
the posterior component weights being defined for count classes with
$P_{\ell,k} > 0$. Consequently the next-answer probability is
$\Pr(\psi(q) = 1\mid\psi(\mathcal S) = y)
= \sum_r w_r(\ell,k)\theta_r$ for $|\mathcal S| = \ell$,
$i \notin \mathcal S$ and $y$ of weight $k$, and averaging its binary
entropy over all count classes gives
%
\begin{equation}
  \gamma_\ell
  = \sum_{\substack{0\le k\le\ell\\ P_{\ell,k}>0}}
  \binom{\ell}{k}P_{\ell,k}\,
  h_2\!\left(\sum_r w_r(\ell,k)\theta_r\right)
  = G_{\ell+1}-G_\ell .
  \label{eq:coingamma}
\end{equation}
%
Zero-probability terms in \eqref{eq:coinGl} use $0\log 0 = 0$, while
\eqref{eq:coingamma} simply omits those classes, whose posterior
weights are undefined and whose contribution vanishes.
Equation~\eqref{eq:coinGl} is valid through $\ell = Q$ and
\eqref{eq:coingamma} through $\ell = Q-1$; apart from that
truncation, neither depends on $n$.

\begin{example}
Take $(\theta_1,\theta_2) = (0.1, 0.9)$ with equal weights, so
$Q = 4$ at $(n,m) = (2,1)$. For a length-$\ell$ string of weight $k$,
equation~\eqref{eq:coinPlk} becomes
$P_{\ell,k} = (9^k + 9^{\ell-k})/(2\cdot 10^\ell)$, so a map's
probability depends only on its Hamming weight, with clearly
separated masses:

\begin{center}
\begin{tabular}{crrr}
\toprule
weight $|\phi_j|$ & 0 & 1 & 2 \\
\midrule
mass of each map & $0.3281$ & $0.0369$ & $0.0081$ \\
multiplicity & 1 & 4 & 6 \\
\bottomrule
\end{tabular}
\end{center}

with weights $3$ and $4$ mirroring weights $1$ and $0$. A weight-one
map is thus more than four times as probable as a weight-two map, and
the normalization $2(0.3281) + 8(0.0369) + 6(0.0081) = 1$ is visible
directly. In the truth-table ordering of the earlier examples:

\begin{center}
\begin{tabular}{llllllll}
\toprule
$j$ & $p_j$ & $j$ & $p_j$ & $j$ & $p_j$ & $j$ & $p_j$ \\
\midrule
0000 & $0.3281$ & 0001 & $0.0369$ & 0010 & $0.0369$ & 0011 & $0.0081$ \\
0100 & $0.0369$ & 0101 & $0.0081$ & 0110 & $0.0081$ & 0111 & $0.0369$ \\
1000 & $0.0369$ & 1001 & $0.0081$ & 1010 & $0.0081$ & 1011 & $0.0369$ \\
1100 & $0.0081$ & 1101 & $0.0369$ & 1110 & $0.0369$ & 1111 & $0.3281$ \\
\bottomrule
\end{tabular}
\end{center}

The mixture mean is $\tfrac12(0.1)+\tfrac12(0.9) = \tfrac12$, so
$G_1 = h_2(1/2) = 1$. If the first answer is $1$, Bayes' rule moves
the component weights from $(\tfrac12,\tfrac12)$ to $(0.1, 0.9)$ and
the next-one probability becomes $0.1(0.1)+0.9(0.9) = 0.82$; if it is
$0$ the weights reverse and that probability is $0.18$. Both branches
therefore have predictive entropy $h_2(0.18) = h_2(0.82)$, so
$G_2 - G_1 = h_2(0.18) = 0.680077$. The remaining block entropies
follow from the same count-by-weight formula: for two answers the
strings $00, 11$ each have probability $0.41$ and $01, 10$ each
$0.09$; for three answers the two constant strings each have
$0.365$ and the other six each $0.045$. Hence
%
\begin{equation}
  (G_0,\ldots,G_4) = (0,\ 1,\ 1.680077,\ 2.269405,\ 2.797921),
  \label{eq:coinworkedG}
\end{equation}
%
with increments $(1,\ 0.680077,\ 0.589327,\ 0.528517)$ and, by the
leverage law \eqref{eq:finitelev}, exact cumulative leverages
$1.95977,\ 1.67930,\ 1.52969,\ 1.42963$. The conditional entropy
density and hyperbola coefficient are
$h_{\rm cond} = h_2(0.1) = 0.468996$ and
$c = (1 - h_2(0.1))/h_2(0.1) = 1.132216$, so the limiting curve is
$L(t) = 1 + 1.132216/t$. The small-$n$ curve lies well below it
because four observations cannot identify the latent coin sharply:
its endpoint is $1.42963$ against the limiting $1 + c = 2.13222$.
\end{example}

\begin{figure}[H]
\centering
\begin{tikzpicture}[baseline=(current bounding box.north)]
\begin{axis}[width=7.4cm, height=5.4cm,
  title={2 components: $\theta = (0.1, 0.9)$},
  title style={font=\small},
  xlabel={$t = \ell/Q$},
  ylabel={$\langle L\rangle$},
  xmin=0.15, xmax=1, ymin=0.95, ymax=9,
  tick label style={font=\scriptsize}, label style={font=\small},
  grid=major, grid style={black!12},
  legend style={font=\tiny, at={(0.97,0.97)}, anchor=north east},
  legend cell align=left]
\dataplot{cn1}{x}{lev}{figures/data/cointwo_n4.dat}
\dataplot{cn2}{x}{lev}{figures/data/cointwo_n6.dat}
\dataplot{cn3}{x}{lev}{figures/data/cointwo_n8.dat}
\dataplot{cn4}{x}{lev}{figures/data/cointwo_n10.dat}
\addplot[clim, dashed, thick, domain=0.15:1, samples=80]
  {1+1.132216/x};
\datalegend{{$n=4$}, {$n=6$}, {$n=8$}, {$n=10$}, {$1+c_1/t$}}
\end{axis}
\end{tikzpicture}\hfill
\begin{tikzpicture}[baseline=(current bounding box.north)]
\begin{axis}[width=7.4cm, height=5.4cm,
  title={10 components, matched $F_1$ and $h_{\rm cond,1}$},
  title style={font=\small},
  xlabel={$t = \ell/Q$},
  xmin=0.15, xmax=1, ymin=0.95, ymax=9,
  tick label style={font=\scriptsize}, label style={font=\small},
  grid=major, grid style={black!12}]
\dataplot{cn1}{x}{lev}{figures/data/cointen_n4.dat}
\dataplot{cn2}{x}{lev}{figures/data/cointen_n6.dat}
\dataplot{cn3}{x}{lev}{figures/data/cointen_n8.dat}
\dataplot{cn4}{x}{lev}{figures/data/cointen_n10.dat}
\addplot[clim, dashed, thick, domain=0.15:1, samples=80]
  {1+1.132216/x};
\end{axis}
\end{tikzpicture}
\caption{Two binary mixtures with the \emph{same} thermodynamic curve
and visibly different finite-size corrections. The right panel uses
ten equally weighted biases $0.005$, $0.04$, $0.1$, $0.2$,
$0.288171$ and their complements, the fifth fitted so that
$\sum h_2 \simeq 5h_2(0.1)$; complement symmetry fixes the mean bias
at $\tfrac12$. Both therefore share $h_0 = 1$ and
$h_{\rm cond} = h_2(0.1)$, hence the same dashed hyperbola, while
their finite curves differ because $I_k$ in \eqref{eq:coindecomp}
depends on the whole collection of weights and separations. Panels
begin at $t = 0.15$ so the common $1/t$ singularity does not dominate
the scale.}
\label{fig:cointdl}
\end{figure}

That sameness is de Finetti's theorem at work: a projectively
consistent exchangeable binary prior keeps only $h_0$ and
$h_{\rm cond}$ in the limit, and two Bernoulli components already
realize every finite $c$ -- two coins are all the richness
exchangeability has.

\subsubsection{What changes when \texorpdfstring{$m>1$}{m>1}}

For $m > 1$ the categorical construction \eqref{eq:coinprior} is the
general model, and the shared-scalar coin prior is a structured
subfamily of it rather than a separate alternative. Use the same
Hamming-weight notation at both scales,
$|a| := \sum_{b=1}^m a_b$ for one $m$-bit answer and
$|\phi_j| := \sum_i |\phi_j(q_i)|$ for the complete table; the signed
bias of Section~\ref{sec:gibbs} is not an independent statistic, being
$B_j = 2|\phi_j| - mQ$. Setting
$\nu_r(a) = \theta_r^{|a|}(1-\theta_r)^{m-|a|}$ draws one component
once for the whole map and makes all $mQ$ table bits independent
Bernoulli variables with bias $\theta_r$ given $Z = r$, so
%
\begin{equation}
  p_j = \sum_r w_r\,\theta_r^{|\phi_j|}(1-\theta_r)^{mQ-|\phi_j|}.
  \label{eq:coinpriorm}
\end{equation}
%
After the shared $Z$ is averaged out, distinct bits are generally
correlated both within one answer and across questions, which is why
the entropy of one $m$-bit answer is not $m\,h_2(\sum_r w_r\theta_r)$.
Writing
%
\begin{equation}
  F_s := -\sum_{k=0}^s\binom sk P_{s,k}\log_2 P_{s,k}
  \label{eq:Fs}
\end{equation}
%
for the entropy of the complete $s$-bit string, and using
$H(\nu_r) = m\,h_2(\theta_r)$, the categorical formulas specialize
exactly to $G_\ell = F_{m\ell}$, $h_0 = F_m$ and
$h_{{\rm cond},m} = m\sum_r w_r h_2(\theta_r)$. The exact finite
leverage plotted below is therefore
%
\begin{equation}
  \langle L_{n,\ell}\rangle
  = \frac{Q F_m - (Q-\ell)\big(F_{m(\ell+1)}-F_{m\ell}\big)}{F_{m\ell}},
  \qquad
  \langle L_{n,Q}\rangle = \frac{Q F_m}{F_{mQ}},
  \label{eq:coinlevm}
\end{equation}
%
and since the latent contributes only subextensive information,
$F_s/s \to \sum_r w_r h_2(\theta_r)$, giving
%
\begin{equation}
  L_m(t) = 1 + \frac{c_m}{t},
  \qquad
  c_m = \frac{F_m - h_{{\rm cond},m}}{h_{{\rm cond},m}} 
  \label{eq:coinhyperbolam}
\end{equation}
%
provided $h_{{\rm cond},m} > 0$. For the mixture above,
$F_1 = 1$, $F_2 = 1.680077$, $h_{{\rm cond},1} = 0.468996$ and
$h_{{\rm cond},2} = 0.937991$, so $c_1 = 1.132216$ and
$c_2 = 0.791144$. The $m = 2$ panel must therefore use $F_2$ for the
initial answer entropy, not $2h_2(1/2) = 2$.

\begin{figure}[H]
\centering
\begin{tikzpicture}[baseline=(current bounding box.north)]
\begin{axis}[width=7.4cm, height=5.4cm,
  title={$m=1$: one bit per answer}, title style={font=\small},
  xlabel={$t = \ell/Q$}, ylabel={$\langle L\rangle$},
  xmin=0.15, xmax=1, ymin=0.95, ymax=9,
  tick label style={font=\scriptsize}, label style={font=\small},
  grid=major, grid style={black!12},
  legend style={font=\tiny, at={(0.97,0.97)}, anchor=north east},
  legend cell align=left]
\dataplot{cn1}{x}{lev}{figures/data/coinm1_n4.dat}
\dataplot{cn2}{x}{lev}{figures/data/coinm1_n6.dat}
\dataplot{cn3}{x}{lev}{figures/data/coinm1_n8.dat}
\dataplot{cn4}{x}{lev}{figures/data/coinm1_n10.dat}
\addplot[clim, dashed, thick, domain=0.15:1, samples=80] {1+1.132216/x};
\datalegend{{$n=4$}, {$n=6$}, {$n=8$}, {$n=10$}, {$1+c_m/t$}}
\end{axis}
\end{tikzpicture}\hfill
\begin{tikzpicture}[baseline=(current bounding box.north)]
\begin{axis}[width=7.4cm, height=5.4cm,
  title={$m=2$: two bits share one bias},
  title style={font=\small},
  xlabel={$t = \ell/Q$},
  xmin=0.15, xmax=1, ymin=0.95, ymax=9,
  tick label style={font=\scriptsize}, label style={font=\small},
  grid=major, grid style={black!12}]
\dataplot{cn1}{x}{lev}{figures/data/coinm2_n4.dat}
\dataplot{cn2}{x}{lev}{figures/data/coinm2_n6.dat}
\dataplot{cn3}{x}{lev}{figures/data/coinm2_n8.dat}
\dataplot{cn4}{x}{lev}{figures/data/coinm2_n10.dat}
\addplot[clim, dashed, thick, domain=0.15:1, samples=80] {1+0.791144/x};
\end{axis}
\end{tikzpicture}
\caption{The same shared-scalar latent at two answer widths. Exact
blocks use $G_\ell = F_{m\ell}$ and the dashes use
\eqref{eq:coinhyperbolam}. The curves approach their hyperbolas at
every fixed $t > 0$; no uniform convergence at the singular origin is
claimed.}
\label{fig:coinm}
\end{figure}

Three superficially similar models are worth separating. With
\emph{one global $Z$ shared by every bit}, the construction used
here, $G_\ell = F_{m\ell}$ and $h_0 = F_m$. With \emph{one
independent global component per output coordinate},
$G_\ell = mF_\ell$, so both initial and residual entropies scale by
$m$ and the coefficient stays $c_1$. With \emph{a fresh component
redrawn at every question}, the answers are i.i.d.\ with distribution
$\nu_{\rm mix}$, so $G_\ell = \ell H(\nu_{\rm mix})$ and $L \equiv 1$:
there is no persistent latent variable to learn.



\subsubsection{Why leverage can have an interior maximum}

The destroyed-entropy density is $h_0 - (1-t)\gamma(t)$, so after
subtracting the received density the cumulative deduced density is
$h_0 - (1-t)\gamma(t) - g(t)$, whose derivative, for differentiable
$\gamma$, is
%
\begin{equation}
  \frac{d}{dt}\Big[h_0-(1-t)\gamma(t)-g(t)\Big]
  = -(1-t)\gamma'(t) \;\ge\; 0 .
  \label{eq:deduceddensity}
\end{equation}
%
Hence $-\gamma'(t)$, weighted by the unasked fraction, specifies when
stored information is released. Since
$L(t) = 1 + \big[h_0-(1-t)\gamma(t)-g(t)\big]/g(t)$,
%
\begin{equation}
  L'(t)
  = \frac{-(1-t)\gamma'(t)g(t)
    - \big[h_0-(1-t)\gamma(t)-g(t)\big]\gamma(t)}{g(t)^2},
  \label{eq:Lprime}
\end{equation}
%
so leverage rises when the current deduction-to-reception rate
exceeds its accumulated average and falls when it drops below it. For
an explicit example mix a fraction $w$ of $(r,k) = (4,2)$ MDS
questions with a fraction $1-w$ of fresh questions, giving
%
\begin{equation}
  \frac{\gamma(t)}m = 1 - 3wt^2 + 2wt^3,
  \qquad
  \frac{g(t)}m = t - wt^3 + \frac w2 t^4,
  \label{eq:peakexample}
\end{equation}
%
and cumulative deduced density
$3wt^2 - 4wt^3 + \tfrac{3w}2 t^4$ per bit. Consequently
$L(t) = 1 + 3wt + O(t^2)$ as $t\downarrow0$, whereas direct
differentiation at the other endpoint gives
%
\begin{equation}
  L'(1) = -\frac w2\,\frac{1-w}{(1-w/2)^2} \;<\; 0
  \qquad (0 < w < 1).
  \label{eq:Lprimeone}
\end{equation}
%
Thus $L'(0^+) = 3w > 0$ while $L'(1) < 0$, and continuity forces an
interior maximum. The peak is produced by a threshold release
followed by continued ordinary reception.

\subsection{An exotic zoo of solvable thermodynamic
priors}\label{sec:zoo}

The following priors are designed to release stored information at
different stages of the run. Their profiles are explicit, so their
leverage follows from \eqref{eq:master}. Full derivations appear in
Appendices~\ref{app:block}--\ref{app:mixtures}.

\subsubsection{Cliques: constant leverage}

For a clique block $\mathcal V \in \mathcal P_n$ of size $r$, draw one
shared value $a_{\mathcal V}\sim\nu_{\mathcal V}$ and copy it to all
questions in $\mathcal V$: the random map answers every question of
the block with that one drawn symbol, $\psi(q) = a_{\mathcal V}$ for
all $q \in \mathcal V$, and the draws are independent across blocks.
Like every zoo member, this is a constructive prior on the $N$ maps.
A map $\phi_j$ restricts to an answer tuple $\phi_j(\mathcal V)$ on
each block, each block law prices that tuple, and independence
multiplies the prices:
%
\begin{equation}
  p_j \;=\; \prod_{\mathcal V\in\mathcal P_n}
  \Pr\big(\psi(\mathcal V) = \phi_j(\mathcal V)\big).
  \label{eq:blockproduct}
\end{equation}
%
For cliques the factor is $\nu_{\mathcal V}(a)$ when $\phi_j$ is
constant on $\mathcal V$ with value $a$, and zero otherwise, so the
support is the $A^{Q/r}$ maps that are constant on every block.

\begin{example}
At $(n,m) = (2,1)$, partition the four questions by their high bit,
$\mathcal V_1 = \{00, 01\}$ and $\mathcal V_2 = \{10, 11\}$, with
$\nu_1 = (0.8, 0.2)$ and $\nu_2 = (0.55, 0.45)$ on the answers
$(0,1)$. A supported map must be constant on both blocks, so it can
depend only on $b_1$: of the sixteen maps of
Table~\ref{tab:sixteen}, only FALSE ($j=0$), $\lnot b_1$ ($j=3$),
$b_1$ ($j=12$) and TRUE ($j=15$) survive, with product weights
\begin{equation}
  p_0 = 0.8\times0.55 = 0.44,\quad
  p_3 = 0.2\times0.55 = 0.11,\quad
  p_{12} = 0.8\times0.45 = 0.36,\quad
  p_{15} = 0.09,
\end{equation}
and $p_j = 0$ for the other twelve. One answer from either block
settles that block's remaining question for free: $L \equiv 2$.
\end{example}

With $h_{\mathcal V} := H(\nu_{\mathcal V})$ we have
$H_{\mathcal V}(0) = 0$ and $H_{\mathcal V}(j) = h_{\mathcal V}$ for
$j \ge 1$, so the increments are
$(h_{\mathcal V}, 0, \ldots, 0)$. For a common block size, letting
$h_{\rm clique}$ denote the question-weighted mean of the
$h_{\mathcal V}$,
%
\begin{equation}
  \gamma(t) = h_{\rm clique}(1-t)^{r-1},
  \qquad
  g(t) = \frac{h_{\rm clique}}{r}\big[1-(1-t)^r\big],
  \qquad
  L(t)\equiv r.
  \label{eq:clique}
\end{equation}
%
The result is exact at finite $n$ for an exact clique partition, or
with only deterministic zero-entropy leftovers. Mixing clique sizes
at a common entropy scale does not interpolate between constants: the
initial leverage is the arithmetic mean of the sizes while the final
leverage is their harmonic mean.

\subsubsection{Parity: deductions near the end}

For a binary block $\mathcal V$ of size $r$, fix a coset label
$\sigma_{\mathcal V}\in\{0,1\}$ and choose its answer tuple uniformly
from that parity coset,
%
\begin{equation}
  \Pr\big(\psi(\mathcal V) = y_{\mathcal V}\big)
  = 2^{-(r-1)}\,\mathbf 1\Big[\textstyle
  \bigoplus_{q\in\mathcal V}y_q = \sigma_{\mathcal V}\Big].
  \label{eq:parityprior}
\end{equation}
%
The same product rule \eqref{eq:blockproduct} makes this
constructive: a map $\phi_j$ is supported precisely when its
restriction to every block lands in that block's coset, and every
supported map carries the same weight
$p_j = \prod_{\mathcal V} 2^{-(r_{\mathcal V}-1)}$.

\begin{example}
At $(n,m) = (2,1)$, let a single parity block cover the whole table,
$\mathcal V = \mathcal Q$ with $r = 4$ and the odd coset
$\sigma_{\mathcal V} = 1$. The support is the eight maps of
Table~\ref{tab:sixteen} whose truth tables have odd weight,
$j \in \{1, 2, 4, 7, 8, 11, 13, 14\}$ -- among them NOR, NAND, AND
and OR -- each with $p_j = 2^{-3} = \tfrac18$. Any three answers admit two
completions of the truth table, of which only one has odd weight,
so the fourth answer is deduced free of charge: the block pays out
exactly at the end.
\end{example}

Every proper subset is uniform, so $\eta(i) = 1$ for
$0\le i\le r-2$ and $\eta(r-1) = 0$, giving
$\gamma(t) = 1-t^{r-1}$, $g(t) = t - t^r/r$ and
%
\begin{equation}
  L(t) = \frac{t + (1-t)t^{r-1}}{t - t^r/r}. 
  \label{eq:parity}
\end{equation}
%
For $r \ge 3$ this rises from $1$ to $r/(r-1)$. The case $r = 2$ is
exceptional: it is an equality or inequality clique, with
$L \equiv 2$. The derivative of the cumulative deduced density,
$-(1-t)\gamma'(t) = (r-1)t^{r-2}(1-t)$, concentrates late in the run.

\subsubsection{Tilted parity}

Uniform parity is not the only solvable law on a parity coset. For
$0 < \theta < 1$ draw i.i.d.\ $\operatorname{Bernoulli}(\theta)$ bits
and condition their parity to be $\sigma$. With
$\varepsilon := 1-2\theta$ and
$\zeta_{r,\sigma} = \tfrac12\big(1+(-1)^\sigma\varepsilon^r\big)$ the
block law is
%
\begin{equation}
  \Pr\nolimits_\sigma\big(\psi(\mathcal V) = y_{\mathcal V}\big)
  = \frac{\theta^{|y_{\mathcal V}|}(1-\theta)^{r-|y_{\mathcal V}|}
    \mathbf 1[\oplus_q y_q = \sigma]}{\zeta_{r,\sigma}},
  \label{eq:tiltedprior}
\end{equation}
%
and if $v$ bits remain unseen with required parity
$\sigma_{\rm rem}$, the next bit is $1$ with probability
%
\begin{equation}
  P_{\rm next}(v,\sigma_{\rm rem})
  = \frac{\theta\big[1-(-1)^{\sigma_{\rm rem}}
    \varepsilon^{\,v-1}\big]}
  {1+(-1)^{\sigma_{\rm rem}}\varepsilon^{\,v}} .
  \label{eq:tiltednext}
\end{equation}
%
A finite average of $h_2(P_{\rm next})$ gives each $\eta(i)$, the
exact weights being derived in Appendix~\ref{app:parity}. The final
member remains determined, so $\eta(r-1) = 0$, while earlier
coefficients now decrease gradually rather than staying exactly flat.
At $r = 4$, $\theta = 0.3$, $\sigma = 1$,
%
\begin{equation}
  \eta = (0.912441,\ 0.898876,\ 0.811317,\ 0),
  \label{eq:tiltedeta}
\end{equation}
%
which gives $L(0^+)\simeq1.0446$ and $L(1)\simeq1.3915$. The
parameter $\varepsilon$ controls how far the profile bends away from
uniform parity.

\subsubsection{MDS blocks: a movable threshold}

Let answers be symbols of $\mathbb F_{A}$ and let a block
$\mathcal V$ of size $r$ be uniform on an $[r,k]_{A}$ MDS code, a
Reed--Solomon realization using $r$ distinct field points and
therefore requiring $A \ge r$. Any $k$ answers determine the whole
block while any $j \le k$ answers are uniform, so
$H_{\rm MDS}(j) = m\min(j,k)$ and $\eta(i) = m$ for $i < k$, zero
afterwards. Writing
$\tau_{r,k}(t) := \Pr[\operatorname{Bin}(r-1,t)\le k-1]$ we obtain
$\gamma(t) = m\,\tau_{r,k}(t)$, and the beta integral gives
%
\begin{equation}
  \rho := \frac kr,
  \qquad
  \int_0^1\tau_{r,k}(x)\,dx = \rho,
  \qquad\text{hence}\qquad
  L(1) = \frac1\rho .
  \label{eq:mdsrate}
\end{equation}
%
If $n\to\infty$ is taken first and then $r\to\infty$ with
$k/r\to\rho$, the profile approaches a step at $t = \rho$, the answer
width having to grow for such blocks to exist. To produce a peak, mix
a code fraction $w$ with a fresh fraction $1-w$:
%
\begin{equation}
  \frac{\gamma(t)}m = w\,\tau_{r,k}(t) + (1-w),
  \qquad
  L(t) = \frac{w\big[1-(1-t)\tau_{r,k}(t)\big] + (1-w)t}
    {w\int_0^t\tau_{r,k}(x)\,dx + (1-w)t} . 
  \label{eq:mdsfresh}
\end{equation}
%
The code releases deductions around its threshold while fresh
questions continue to add received information afterwards, making the
cumulative curve decline past its maximum.

\subsubsection{Local noise and full support}

The hard clique, parity and code priors assign probability zero to
many maps. For $0\le\delta\le1$ a general local lapse replaces the
hard law $P_{\mathcal V}$ on a block by
$P_{\mathcal V,\delta} = (1-\delta)P_{\mathcal V}
+ \delta A^{-|\mathcal V|}$, so that term by term the map prior
becomes
%
\begin{equation}
  p_{\delta,j}
  = \prod_{\mathcal V\in\mathcal P_n}
  \Big[(1-\delta_{\mathcal V,n})
    P_{\mathcal V}\big(\phi_j(\mathcal V)\big)
    + \delta_{\mathcal V,n}A^{-|\mathcal V|}\Big] . 
  \label{eq:lapse}
\end{equation}
%
The map prior has full support whenever every
$\delta_{\mathcal V,n} > 0$, and the lapse is local: blocks remain
independent, so only their entropy coefficients change. For fixed
block sizes entropy is continuous on the finite probability simplex,
so $\eta_{\mathcal V,\delta}(i)\to\eta_{\mathcal V,0}(i)$ as
$\delta\to0$, and if
$\max_{\mathcal V}\delta_{\mathcal V,n}\to0$ the hard and noisy
priors share a thermodynamic profile. More generally it is enough
that the question-weighted coefficient perturbation vanish,
%
\begin{equation}
  \sum_{\mathcal V\in\mathcal P_n}
  \frac{|\mathcal V|}{Q}
  \max_i\big|\eta_{\mathcal V,\delta_{\mathcal V,n}}(i)
  - \eta_{\mathcal V,0}(i)\big| \;\longrightarrow\; 0 .
  \label{eq:noisecondition}
\end{equation}
%
At fixed $\delta > 0$ the noise does \emph{not} vanish merely because
$n\to\infty$; it produces a solvable but displaced profile. If block
sizes also grow, $\delta_n\to0$ alone may be insufficient: for noisy
parity blocks of size $r(n)$ one needs roughly
$r(n)\delta_n\to0$. Two coefficient changes are especially simple.
Uniform parity followed by independent bit flips at rate $\delta$
keeps $\eta(i) = 1$ for $i\le r-2$ and lifts only the cliff,
%
\begin{equation}
  \eta(r-1)
  = h_2\!\left(\frac{1-(1-2\delta)^r}{2}\right),
  \label{eq:noisyparityeta}
\end{equation}
%
while an MDS block lapsed to the uniform block distribution keeps
$H_\delta(j) = jm$ below threshold and turns the zero coefficients
above it into a positive tail, with the exact entropy given in
Appendix~\ref{app:mds}. Merely giving the shared clique value full
alphabet support is not enough, since the block would still contain
only constant strings: independent output noise or a lapse is
required for full map support.

\subsubsection{Arbitrary mixtures and the cocktail}

Suppose component type $\kappa$ occupies a fraction $w_\kappa$ of
questions and has initial marginal entropy $h_{0,\kappa}$, profile
$\gamma_\kappa$, integral $g_\kappa$ and leverage $L_\kappa$.
Independence across disjoint question sets gives
$h_0 = \sum_\kappa w_\kappa h_{0,\kappa}$,
$\gamma = \sum_\kappa w_\kappa\gamma_\kappa$ and
$g = \sum_\kappa w_\kappa g_\kappa$, and therefore
%
\begin{equation}
  L(t)
  = \frac{\sum_\kappa w_\kappa g_\kappa(t)L_\kappa(t)}
         {\sum_\kappa w_\kappa g_\kappa(t)} . 
  \label{eq:mixture}
\end{equation}
%
Leverage is a received-entropy-weighted average of component
leverages, not a simple question-fraction average, and
\eqref{eq:mixture} is a symbolic solution for any mixture whose
component profiles are known.

As a worked cocktail take $m = 4$, with one quarter of the questions
following the two-coin mixture $\theta\in\{0.2,0.7\}$ with weights
$(\tfrac14,\tfrac34)$ and one scalar latent bias shared by all its
bits; one half belonging to lapsed $(r,k) = (16,4)$ Reed--Solomon
blocks with lapse $\delta = 0.1$; and one quarter fresh with
independent bit bias $0.3$. Normalizing entropy per output bit and
writing $h_{\rm cocktail} := h_0/m$,
$\gamma_{\rm cocktail} := \gamma/m$ and
$g_{\rm cocktail} := g/m$, the coin component contributes
$F_4/4 = 0.94828746$ initially and
$h_{\rm cond,coin} = 0.84145020$ afterwards, so
%
\begin{equation}
  h_{\rm cocktail}
  = \tfrac14(0.94828746) + \tfrac12 + \tfrac14 h_2(0.3)
  = 0.95739459,
  \label{eq:cocktailh0}
\end{equation}
%
while for $t > 0$, with $\gamma_{\rm lapse}$ the lapsed-code profile
per bit,
%
\begin{equation}
  \gamma_{\rm cocktail}(t)
  = \tfrac14 h_{\rm cond,coin} + \tfrac12\gamma_{\rm lapse}(t)
  + \tfrac14 h_2(0.3),
  \qquad
  \int_0^1\gamma_{\rm lapse} = 0.33232806 .
  \label{eq:cocktailgamma}
\end{equation}
%
Substituting into \eqref{eq:master} gives
$L = \big[h_{\rm cocktail}-(1-t)\gamma_{\rm cocktail}\big]
/g_{\rm cocktail}$, whose initial divergence has coefficient
$\lim_{t\downarrow0}tL(t)
= \big[h_{\rm cocktail}-\gamma_{\rm cocktail}(0^+)\big]
/\gamma_{\rm cocktail}(0^+)\simeq0.02870$, a dip of approximately
$1.499$ at $t\simeq0.0836$, a peak of approximately $2.114$ at
$t\simeq0.350$, and $L(1)\simeq1.60408$. The initial hyperbolic
decline comes from the global coin latent, the middle peak from the
code threshold, and the final decline from fresh and lapsed entropy
received after the code deductions have mostly been released.

\begin{figure}[H]
\centering
\begin{tikzpicture}[baseline=(current bounding box.north)]
\begin{axis}[width=4.05cm, height=4.6cm,
  title={cliques}, title style={font=\scriptsize},
  xlabel={$t$}, ylabel={$\langle L\rangle$},
  xmin=0, xmax=1, ymin=0.9, ymax=4.6,
  tick label style={font=\tiny}, label style={font=\scriptsize},
  grid=major, grid style={black!12}]
\dataplot{cn3}{x}{lev}{figures/data/zoo_clique_r2.dat}
\dataplot{cn3}{x}{lev}{figures/data/zoo_clique_r4.dat}
\addplot[clim, dashed, domain=0:1] {2};
\addplot[clim, dashed, domain=0:1] {4};
\end{axis}
\end{tikzpicture}\hfill
\begin{tikzpicture}[baseline=(current bounding box.north)]
\begin{axis}[width=4.05cm, height=4.6cm,
  title={parity, $r=4$}, title style={font=\scriptsize},
  xlabel={$t$}, xmin=0, xmax=1, ymin=0.97, ymax=1.4,
  tick label style={font=\tiny}, label style={font=\scriptsize},
  grid=major, grid style={black!12}]
\dataplot{cn1}{x}{lev}{figures/data/zoo_parity_n6.dat}
\dataplot{cn3}{x}{lev}{figures/data/zoo_parity_n8.dat}
\dataplot{cn4}{x}{lev}{figures/data/zoo_parity_n10.dat}
\addplot[clim, dashed, domain=0.02:1, samples=80]
  {(x+(1-x)*x^3)/(x-x^4/4)};
\end{axis}
\end{tikzpicture}\hfill
\begin{tikzpicture}[baseline=(current bounding box.north)]
\begin{axis}[width=4.05cm, height=4.6cm,
  title={MDS $+$ fresh}, title style={font=\scriptsize},
  xlabel={$t$}, xmin=0, xmax=1, ymin=0.9, ymax=3.2,
  tick label style={font=\tiny}, label style={font=\scriptsize},
  grid=major, grid style={black!12}]
\dataplot{cn1}{x}{lev}{figures/data/zoo_mds_n6.dat}
\dataplot{cn3}{x}{lev}{figures/data/zoo_mds_n8.dat}
\dataplot{cn4}{x}{lev}{figures/data/zoo_mds_n10.dat}
\dataplot{clim, dashed}{x}{lev}{figures/data/zoo_mds_limit.dat}
\end{axis}
\end{tikzpicture}\hfill
\begin{tikzpicture}[baseline=(current bounding box.north)]
\begin{axis}[width=4.05cm, height=4.6cm,
  title={cocktail}, title style={font=\scriptsize},
  xlabel={$t$}, xmin=0, xmax=1, ymin=0.9, ymax=4,
  tick label style={font=\tiny}, label style={font=\scriptsize},
  grid=major, grid style={black!12}]
\dataplot{cn1}{x}{lev}{figures/data/zoo_cocktail_n6.dat}
\dataplot{cn3}{x}{lev}{figures/data/zoo_cocktail_n8.dat}
\dataplot{cn4}{x}{lev}{figures/data/zoo_cocktail_n10.dat}
\dataplot{clim, dashed}{x}{lev}{figures/data/zoo_cocktail_limit.dat}
\end{axis}
\end{tikzpicture}
\caption{Finite learning curves at $n = 6, 8, 10$ (light to dark)
against their dashed thermodynamic limits. Cliques of size $2$ and
$4$ are exactly flat at every finite size; parity rises to
$r/(r-1)$; the diluted code peaks near its threshold; and the
cocktail declines, dips, peaks and declines again. In the cocktail
panel the dashed curve uses the shared-latent initial entropy $F_4/4$
and equations \eqref{eq:cocktailh0}--\eqref{eq:cocktailgamma}.}
\label{fig:zoo}
\end{figure}

\subsubsection{The zoo side by side}

\begin{center}
\begin{tabular}{@{}llll@{}}
\toprule
member & microscopic rule & limiting profile & leverage shape \\
\midrule
coin mixture & a global latent selects an i.i.d.\ component
  & constant on $t>0$, jump at $0$ & $1 + c/t$ \\
clique & one value copied through each block
  & $h_{\rm clique}(1-t)^{r-1}$ & exactly $r$ \\
parity & one final parity constraint
  & $m(1-t^{r-1})$ & rising for $r\ge3$ \\
tilted parity & biased product law conditioned on parity
  & \eqref{eq:bernsteinprofile} with \eqref{eq:tiltedeta}
  & tilt-dependent \\
MDS & any $k$ answers determine an $r$-block
  & $m\,\tau_{r,k}(t)$ & threshold rise \\
MDS $+$ fresh & threshold block plus ordinary sites
  & $m[w\tau_{r,k}+(1-w)]$ & interior peak \\
cocktail & coin, lapsed code and fresh
  & weighted sum & decline, dip, peak \\
\bottomrule
\end{tabular}
\end{center}

The shape of a well-measured learning curve can therefore provide
evidence about the \textbf{timing} of usable correlations: global and
front-loaded, first-touch, last-touch, or threshold-like. It does not
identify a unique microscopic prior. Different priors can share the
same entropy profile, and an actual trained model may be
misspecified, non-Bayesian, or imperfectly calibrated. The central
scale distinction is now explicit: the $n\to\infty$ curve sees only
extensive-time learning, and any structure learned in $O(1)$ or
$o(Q)$ questions collapses into a singularity or jump at $t = 0$. For
Bayesian priors this means that a finite-dimensional global latent may
profoundly affect early learning yet disappear from the interior
profile. To obtain nontrivial behaviour throughout $0 < t < 1$, the
prior must distribute uncertainty and correlations across $\Theta(Q)$
local degrees of freedom. Macroscopic leverage classifies when stored
information becomes available, not every detail of how the prior
represents it.

\appendix

Appendix~\ref{app:spike} continues the spike-prior example of
Section~\ref{sec:spike}.
Appendices~\ref{app:block}--\ref{app:mixtures} derive the block
formulas of Chapter~\ref{sec:tdl} in a common order: specify the
probability law inside one block; compute its subset entropies
$H_\kappa(j)$; take differences
$\eta_\kappa(i) = H_\kappa(i+1)-H_\kappa(i)$; insert those
differences into the Bernstein profile; integrate the profile and
substitute it into the leverage law.
Appendix~\ref{app:inverse} discusses the inverse problem and
Appendix~\ref{app:endpoints} collects endpoint and boundary-layer
consequences of the general limit.


\section{The general block calculation}\label{app:block}

\subsection{Exact mean block entropy at finite \texorpdfstring{$n$}{n}}

Let $\mathcal V \in \mathcal P_n$ be a block and define the entropy
averaged over all $j$-subsets of it,
%
\begin{equation}
  H_{\mathcal V}(j)
  := \binom{|\mathcal V|}{j}^{-1}
  \sum_{\substack{\mathcal J\subseteq\mathcal V\\ |\mathcal J| = j}}
  H\big(\psi(\mathcal J)\big),
  \label{eq:appblockH}
\end{equation}
%
every term of which is equal when the block is entropy homogeneous.
If $\mathcal S$ is a uniformly random $\ell$-subset of $\mathcal Q$
then its intersection count with $\mathcal V$ is hypergeometric,
%
\begin{equation}
  \Pr(|\mathcal S\cap\mathcal V| = j)
  = \frac{\binom{|\mathcal V|}{j}\binom{Q-|\mathcal V|}{\ell-j}}
         {\binom{Q}{\ell}} .
  \label{eq:apphyper}
\end{equation}
%
Different blocks have dependent intersection counts, since those
counts sum to $\ell$, but entropy adds across independent blocks and
expectation is linear, so
%
\begin{equation}
  G_{n,\ell}
  = \sum_{\mathcal V\in\mathcal P_n}\sum_{j=0}^{|\mathcal V|}
  \frac{\binom{|\mathcal V|}{j}\binom{Q-|\mathcal V|}{\ell-j}}
       {\binom{Q}{\ell}}\,H_{\mathcal V}(j) 
  \label{eq:appGnl}
\end{equation}
%
exactly, with no large-$n$ approximation anywhere.

\subsection{Exact conditional-entropy increments}

Choose a fresh question $q$ uniformly and then $\ell$ of the
remaining $Q-1$ questions. Conditional on $q$ belonging to a
type-$\kappa$ block of size $r_\kappa$, the probability of seeing
exactly $i$ already-asked partners is
$\binom{r_\kappa-1}{i}\binom{Q-r_\kappa}{\ell-i}\big/
\binom{Q-1}{\ell}$, and only those partners affect the fresh answer
because blocks are independent. Choosing an $i$-subset and then a
fresh coordinate uniformly shows that
$\eta_\kappa(i) = H_\kappa(i+1)-H_\kappa(i)$ is the corresponding
mean conditional entropy even without entropy homogeneity, so for
$0\le\ell<Q$,
%
\begin{equation}
  \gamma_{n,\ell}
  = G_{n,\ell+1}-G_{n,\ell}
  = \sum_\kappa w_{\kappa,n}\sum_{i=0}^{r_\kappa-1}
  \frac{\binom{r_\kappa-1}{i}\binom{Q-r_\kappa}{\ell-i}}
       {\binom{Q-1}{\ell}}\,\eta_\kappa(i) . 
  \label{eq:appgamma}
\end{equation}

\subsection{Hypergeometric to binomial}

Fix $r$ and $i$, set $\ell = \lfloor tQ\rfloor$, and expose a
particular set of $i$ partners that are selected and $r-1-i$ that are
not. Its probability is a product of $r-1$ factors of the form
$(\ell+O(r))/(Q+O(r))$ or $(Q-\ell+O(r))/(Q+O(r))$; each selected
factor tends to $t$ and each unselected factor to $1-t$. There are
$\binom{r-1}{i}$ such choices, so
$\Pr(\text{exactly } i \text{ partners})\to
\binom{r-1}{i}t^i(1-t)^{r-1-i}$, with total-variation error
$O(r^2/Q)$ at fixed $r$. Thus \eqref{eq:appgamma} tends to
\eqref{eq:bernsteinprofile}.

\subsection{Integrating a Bernstein profile}

Direct integration of one basis function gives
%
\begin{equation}
  \int_0^t\beta_{r-1,i}(x)\,dx
  = \binom{r-1}{i}\int_0^t x^i(1-x)^{r-1-i}\,dx,
  \qquad
  \int_0^1\beta_{r-1,i}(x)\,dx = \frac1r,
  \label{eq:appbetaint}
\end{equation}
%
so that
$g(t) = \sum_\kappa w_\kappa\sum_i\eta_\kappa(i)\binom{r_\kappa-1}{i}
\int_0^t x^i(1-x)^{r_\kappa-1-i}dx$ and, at the end of the run,
%
\begin{equation}
  g(1)
  = \sum_\kappa\frac{w_\kappa}{r_\kappa}\sum_i\eta_\kappa(i)
  = \sum_\kappa\frac{w_\kappa}{r_\kappa}H_\kappa(r_\kappa),
  \label{eq:appgone}
\end{equation}
%
exactly the prior entropy per question: there are
$w_\kappa Q/r_\kappa$ independent type-$\kappa$ blocks, each carrying
$H_\kappa(r_\kappa)$ bits.

\subsection{MDS mixtures as Bernstein coefficient synthesizers}

Fix a common block length $r$ over $\mathbb F_{A}$ with $A \ge r$, so
that $[r,k]_{A}$ Reed--Solomon codes exist for $k = 1,\ldots,r$, and
normalize entropies by $m$. A dimension-$k$ block contributes
coefficient $1$ at indices $i < k$ and $0$ afterwards, so if
dimension $k$ is used with question fraction $w_k$ the mixed
normalized coefficients are
$\eta_{\rm target}(i) = \sum_{k=i+1}^{r}w_k$. Conversely, take any
nonincreasing sequence
$1\ge\eta_{\rm target}(0)\ge\cdots\ge\eta_{\rm target}(r-1)\ge0$, set
$\eta_{\rm target}(r) = 0$, and choose
%
\begin{equation}
  w_k = \eta_{\rm target}(k-1)-\eta_{\rm target}(k)
  \qquad (1\le k\le r),
  \label{eq:appsynth}
\end{equation}
%
which is nonnegative and reproduces the sequence, any slack
$1-\eta_{\rm target}(0)$ being assigned to frozen zero-entropy
blocks. Bernstein's approximation theorem then approximates a
continuous nonincreasing normalized target profile by taking
$\eta_{\rm target}(i) = \gamma_{\rm target}(i/(r-1))$. This
establishes a profile-synthesis result only under the accompanying
realization conditions, $r^2/Q\to0$ along a diagonal limit, $A\ge r$,
and leftovers of vanishing density; it is not a completeness theorem
at fixed $m$.

\section{Cliques}\label{app:cliques}

\subsection{Block law and entropy vector}

Let $\mathcal V\in\mathcal P_n$ have size $r$. Draw an $m$-bit shared
value $a_{\mathcal V}\sim\nu_{\mathcal V}$ and set $\psi(q) =
a_{\mathcal V}$ for every $q\in\mathcal V$, so that
$\Pr(\psi(\mathcal V) = y_{\mathcal V})
= \nu_{\mathcal V}(y_{q_1})\mathbf 1[y_{q_1} = \cdots = y_{q_r}]$.
Observing no coordinate has zero entropy; observing at least one
reveals $a_{\mathcal V}$, after which further coordinates add
nothing:
%
\begin{equation}
  H_{\mathcal V}(0) = 0,
  \qquad
  H_{\mathcal V}(j) = h_{\mathcal V}\ (1\le j\le r),
  \qquad
  \eta_{\mathcal V} = (h_{\mathcal V}, 0, \ldots, 0).
  \label{eq:appcliqueeta}
\end{equation}

\subsection{Profile, integral, and leverage}

Only the $i = 0$ Bernstein term survives, so
$\gamma(t) = h_{\rm clique}(1-t)^{r-1}$ and, substituting
$u = 1-x$, $g(t) = \tfrac{h_{\rm clique}}{r}[1-(1-t)^r]$. The
initial marginal entropy is $h_0 = h_{\rm clique}$ and the destroyed
density is
%
\begin{equation}
  h_0-(1-t)\gamma(t)
  = h_{\rm clique}\big[1-(1-t)^r\big]
  = r\,g(t),
  \qquad\text{hence}\qquad
  L(t) = r .
  \label{eq:appcliqueL}
\end{equation}
%
For example, one four-question clique with a fair binary value has
prior mass $\tfrac12$ on $0000$ and $\tfrac12$ on $1111$; its profile
is $(1-t)^3$, its received density $[1-(1-t)^4]/4$, and its leverage
$4$ at every stage.

\subsection{Why the result is exact at finite size}

A block is learned if and only if it has been touched at least once,
and its probability of remaining untouched after a random
$\ell$-subset is $\binom{Q-r}{\ell}\big/\binom{Q}{\ell}$. For $Q/r$
equal-entropy blocks the expected received information is
$\tfrac{Q}{r}h_{\rm clique}\big[1-\binom{Q-r}{\ell}/\binom{Q}{\ell}\big]$
while the expected destroyed table entropy is
$Q\,h_{\rm clique}\big[1-\binom{Q-r}{\ell}/\binom{Q}{\ell}\big]$.
Their ratio is exactly $r$ at every non-null stage.

\subsection{Mixing clique sizes}

Suppose a question fraction $w_\kappa$ lies in size-$r_\kappa$
cliques with a common entropy scale. Substituting each type's
destroyed factor $1-(1-t)^{r_\kappa}$ and received factor
$r_\kappa^{-1}[1-(1-t)^{r_\kappa}]$ gives
%
\begin{equation}
  L(t)
  = \frac{\sum_\kappa w_\kappa\big[1-(1-t)^{r_\kappa}\big]}
         {\sum_\kappa (w_\kappa/r_\kappa)
           \big[1-(1-t)^{r_\kappa}\big]},
  \qquad
  L(0^+) = \sum_\kappa w_\kappa r_\kappa,
  \qquad
  L(1) = \Big(\sum_\kappa\frac{w_\kappa}{r_\kappa}\Big)^{-1},
  \label{eq:appcliquemix}
\end{equation}
%
the arithmetic and harmonic means of the sizes. A mixture of flat
clique curves is therefore generally a decreasing curve, not another
constant. If type-$\kappa$ cliques instead have mean entropy
$h_\kappa$, the entropy-weighted endpoints are
$L(0^+) = \sum w_\kappa h_\kappa r_\kappa/\sum w_\kappa h_\kappa$ and
$L(1) = \sum w_\kappa h_\kappa/\sum w_\kappa h_\kappa/r_\kappa$.

\subsection{Noisy cliques}

For binary answers draw the shared bit $Z\sim\nu$ and let
$\psi(q_j) = Z\oplus E_j$ with
$E_j$ i.i.d.\ $\operatorname{Bernoulli}(\delta)$. Writing
$\mathcal J_i := \{q_1,\ldots,q_i\}$, the posterior on the shared
value after observing $y_{\mathcal J_i}$ is
%
\begin{equation}
  P_i(z\mid y_{\mathcal J_i})
  \;\propto\;
  \nu(z)\prod_{j=1}^{i}
  \big[(1-\delta)\mathbf 1(y_{q_j} = z)
   + \delta\mathbf 1(y_{q_j}\ne z)\big],
  \label{eq:appnoisyclique}
\end{equation}
%
so the next-one probability is
$\delta + (1-2\delta)P_i(1\mid y_{\mathcal J_i})$ and
%
\begin{equation}
  \eta_\delta(i)
  = \sum_{y_{\mathcal J_i}\in\{0,1\}^i}
  \Pr\big(\psi(\mathcal J_i) = y_{\mathcal J_i}\big)\,
  h_2\!\big(\delta + (1-2\delta)P_i(1\mid y_{\mathcal J_i})\big).
  \label{eq:appnoisycliqueeta}
\end{equation}
%
For every $0<\delta<1$ all binary strings have positive probability,
and as $\delta\to0$ this tends to \eqref{eq:appcliqueeta} term by
term. At fixed nontrivial noise the exact $L\equiv r>1$ law is lost,
and at $\delta = \tfrac12$ the output is i.i.d.\ fair and the
different flat law $L\equiv1$ results.

\section{Uniform, tilted, and noisy parity}\label{app:parity}

\subsection{Uniform parity}

Fix $\sigma\in\{0,1\}$ and choose the binary answers in
$\mathcal V$ uniformly subject to
$\bigoplus_{q\in\mathcal V}\psi(q) = \sigma$. There are $2^{r-1}$
admissible strings; fixing any $j\le r-1$ coordinates at any values
leaves the remaining $r-j$ obeying one parity equation and therefore
$2^{r-j-1}$ completions, independent of the fixed values. Every
proper subset is consequently uniform,
%
\begin{equation}
  H_{\mathcal V}(j) = j\ \ (0\le j\le r-1),
  \qquad
  H_{\mathcal V}(r) = r-1,
  \qquad
  \eta(i) =
  \begin{cases}
    1, & 0\le i\le r-2,\\
    0, & i = r-1 .
  \end{cases}
  \label{eq:apparityeta}
\end{equation}
%
For $m$-bit answers the same formulas acquire a factor $m$ if parity
is imposed independently on each bit, or as one additive equation
over $\mathbb F_{A}$.

\subsection{Profile and leverage}

The Bernstein basis sums to one, so removing only the final term
gives $\gamma(t) = \sum_{i=0}^{r-2}\beta_{r-1,i}(t) = 1-t^{r-1}$ and
$g(t) = t-t^r/r$. The destroyed density is
$1-(1-t)(1-t^{r-1}) = t+(1-t)t^{r-1}$, so $L$ is
\eqref{eq:parity}, and for $r\ge3$ direct differentiation yields
%
\begin{equation}
  L'(t)
  = \frac{t^{r-1}(1-t)
    \big[r(r-2)-\sum_{j=1}^{r-2}t^j\big]}
    {r\big(t-t^r/r\big)^2} \;\ge\; 0,
  \qquad
  L(0^+) = 1,
  \qquad
  L(1) = \frac r{r-1},
  \label{eq:apparityLprime}
\end{equation}
%
with $L\equiv2$ at $r = 2$. For the smallest genuinely different
example take one odd-parity block with $r = 4$, uniform on the eight
odd-weight strings, so that $\gamma(t) = 1-t^3$, $g(t) = t-t^4/4$.
The derivative of the cumulative deduced density,
$-(1-t)\gamma'(t) = (r-1)t^{r-2}(1-t)$, is proportional to a
$\operatorname{Beta}(r-1,2)$ density; it is also the density of the
stage at which a block's penultimate queried member determines its
final member, divided by $r$ questions per block.

\subsection{Tilted parity: exact subset probabilities}

For $0<\theta<1$ start with i.i.d.\
$\operatorname{Bernoulli}(\theta)$ bits and condition on the parity.
With $\varepsilon = 1-2\theta$, the probability that $v$ i.i.d.\ bits
have parity $\sigma_{\rm rem}$ is
$\zeta_{v,\sigma_{\rm rem}}
= \tfrac12\big(1+(-1)^{\sigma_{\rm rem}}\varepsilon^v\big)$, and the
normalized tilted-coset law is \eqref{eq:tiltedprior}. Take a
particular realized tuple $y_{\mathcal J}$ on
$\mathcal J\subset\mathcal V$ with $|\mathcal J| = j$ and Hamming
weight $z$; its unobserved coordinates must have parity
$\sigma\oplus(z\bmod2)$, and summing them out gives
%
\begin{equation}
  \Pr\nolimits_\sigma\big(\psi(\mathcal J) = y_{\mathcal J}\big)
  = \theta^z(1-\theta)^{j-z}\,
  \frac{1+(-1)^{\sigma\oplus(z\bmod2)}\varepsilon^{\,r-j}}
       {1+(-1)^\sigma\varepsilon^{\,r}}
  \;=:\; \chi_\sigma(j,z).
  \label{eq:apptiltedsubset}
\end{equation}
%
The subset entropy is then the finite sum
$H_\sigma(j) = -\sum_{z=0}^j\binom jz\chi_\sigma(j,z)
\log_2\chi_\sigma(j,z)$, and the coefficients are exactly
$\eta(i) = H_\sigma(i+1)-H_\sigma(i)$.

\subsection{Tilted parity: predictive form}

The same result can be organized as a conditional prediction. Suppose
$|\mathcal J| = i$ questions have been seen, so $v := r-i$ bits
remain; if their required parity is $\sigma_{\rm rem}$ then
$P_{\rm next}(v,\sigma_{\rm rem})$ is \eqref{eq:tiltednext}. Letting
$\sigma_{\rm seen}$ be the parity of the $i$ seen bits, conditional
on total parity $\sigma$,
%
\begin{equation}
  \Pr(\sigma_{\rm seen} = z\mid\sigma)
  = \frac{\big[1+(-1)^z\varepsilon^{\,i}\big]
    \big[1+(-1)^{\sigma\oplus z}\varepsilon^{\,v}\big]}
    {2\big[1+(-1)^\sigma\varepsilon^{\,r}\big]},
  \qquad
  \eta(i) = \sum_{z=0}^1\Pr(\sigma_{\rm seen} = z\mid\sigma)\,
  h_2\big(P_{\rm next}(v,\sigma\oplus z)\big). 
  \label{eq:apptiltedeta}
\end{equation}
%
When $\theta = \tfrac12$ we have $\varepsilon = 0$ and this reduces
to \eqref{eq:apparityeta}; when $v = 1$ the remaining bit is fixed and
$\eta(r-1) = 0$. At $r = 4$, $\theta = 0.3$, $\sigma = 1$, evaluation
gives \eqref{eq:tiltedeta}.

\subsection{Uniform parity followed by bit-flip noise}

Return to the uniform coset and flip every bit independently with
probability $\delta$. The parity changes precisely when an odd number
of flips occurs, with
$P_{\rm odd}(r,\delta) = \tfrac12\big[1-(1-2\delta)^r\big]$, so
convolution with the flip channel produces a mixture of the two
uniform parity cosets with weights $1-P_{\rm odd}$ and
$P_{\rm odd}$. Any $r-1$ coordinates remain uniform, so
$\eta_\delta(i) = 1$ for $i\le r-2$, while given $r-1$ coordinates
the uncertainty in the final bit is exactly the uncertainty in the
resulting parity label:
%
\begin{equation}
  \gamma_\delta(t)
  = 1-\big[1-h_2(P_{\rm odd}(r,\delta))\big]t^{r-1}. 
  \label{eq:appnoisyparity}
\end{equation}
%
For fixed $r$ any $\delta_n\to0$ recovers the hard profile; if the
block size grows one instead needs
$P_{\rm odd}(r(n),\delta_n)\to0$, which for small noise requires
roughly $r(n)\delta_n\to0$. These formulas rely on a uniform starting
coset: tilted parity followed by noise remains solvable by finite
sums but is not uniform within each resulting coset.

\section{MDS and Reed--Solomon blocks}\label{app:mds}

\subsection{Block construction}

Work over $\mathbb F_{A}$ and choose $r$ distinct evaluation points
$\alpha_1,\ldots,\alpha_r$, which requires $r \le A$. Draw a
polynomial $f(z) = d_0 + d_1z + \cdots + d_{k-1}z^{k-1}$ uniformly
from the $A^k$ polynomials of degree below $k$ and set
$\psi(q_j) = f(\alpha_j)$, a uniform Reed--Solomon codeword. Two
elementary interpolation facts determine all entropies: values at any
$k$ distinct points determine $f$ and hence all $r$ answers; and
values at any $j\le k$ points are uniform on $\mathbb F_{A}^j$, since
prescribing those $j$ values, choosing values at another $k-j$ points
freely and interpolating shows that exactly $A^{k-j}$ polynomials
realize every prescription. Therefore
%
\begin{equation}
  H_{\rm MDS}(j) = m\min(j,k),
  \qquad
  \eta(i) =
  \begin{cases}
    m, & 0\le i\le k-1,\\
    0, & k\le i\le r-1 .
  \end{cases}
  \label{eq:appmdseta}
\end{equation}
%
The same entropy law holds for any uniform coset of an $[r,k]_{A}$
MDS code. Repetition ($k = 1$) and single-parity-check ($k = r-1$)
codes give the clique and parity entropy laws over any alphabet,
although for small fields and large $r$ they are not obtained from
$r$ distinct Reed--Solomon evaluation points.

\subsection{Profile and code rate}

Define the binomial lower-tail polynomial
$\tau_{r,k}(t) := \sum_{i=0}^{k-1}\binom{r-1}{i}t^i(1-t)^{r-1-i}$.
Inserting \eqref{eq:appmdseta} into the Bernstein law gives
$\gamma(t) = m\,\tau_{r,k}(t)$, whose integral is
%
\begin{equation}
  \int_0^1\tau_{r,k}(x)\,dx
  = \sum_{i=0}^{k-1}\binom{r-1}{i}\frac{i!\,(r-1-i)!}{r!}
  = \sum_{i=0}^{k-1}\frac1r
  = \rho,
  \qquad \rho := \frac kr,
  \label{eq:appmdsrate}
\end{equation}
%
the same result following from dividing the $km$ bits per block by
$r$ questions. For a pure code prior
$g(t) = m\int_0^t\tau_{r,k}$ and
%
\begin{equation}
  L_{\rm code}(t)
  = \frac{1-(1-t)\tau_{r,k}(t)}{\int_0^t\tau_{r,k}(x)\,dx},
  \qquad
  L_{\rm code}(1) = \frac1\rho = \frac rk . 
  \label{eq:appmdsL}
\end{equation}
%
For $A = 4$, $r = 4$ and $k = 2$, a random affine polynomial over
$\mathbb F_4$ gives $16$ codewords, any two answers determining the
other two, and
$\tau_{4,2}(t) = (1-t)^3+3t(1-t)^2 = 1-3t^2+2t^3$, the threshold
component of the explicit non-monotonic example
\eqref{eq:peakexample}.

\subsection{The secondary large-block limit}

Now take $r\to\infty$ with $k/r\to\rho$, \emph{after} taking
$n\to\infty$ at fixed $r$. By the law of large numbers a
$\operatorname{Bin}(r-1,t)$ count divided by $r-1$ converges in
probability to $t$, so away from $t = \rho$,
$\tau_{r,k}(t)\to\mathbf 1[t<\rho]$ and hence
%
\begin{equation}
  L_{\rm code}(t)\longrightarrow
  \begin{cases}
    1, & t<\rho,\\
    1/\rho, & t>\rho .
  \end{cases}
  \label{eq:appmdsstep}
\end{equation}
%
At finite $r$ the transition is smooth and $\gamma(t)>0$ for every
$t<1$; the jump appears only in the secondary limit, and a
Reed--Solomon realization of it requires $m\ge\log_2 r$ to grow.

\subsection{Dilution with fresh questions}

Let a fraction $w$ of questions lie in code blocks and $1-w$ be
independent uniform answers, so that
$\gamma(t)/m = w\tau_{r,k}(t)+(1-w)$ and
$g(t)/m = w\int_0^t\tau_{r,k}+(1-w)t$; since the initial marginal
entropy is $m$, the leverage is \eqref{eq:mdsfresh}. For
$1\le k<r$ differentiating the binomial tail gives
$\tau_{r,k}'(t) = -(r-1)\binom{r-2}{k-1}t^{k-1}(1-t)^{r-1-k}$, so
%
\begin{equation}
  -(1-t)\gamma'(t)
  = w\,m\,(r-1)\binom{r-2}{k-1}t^{k-1}(1-t)^{r-k} 
  \label{eq:appmdsrelease}
\end{equation}
%
is proportional to a $\operatorname{Beta}(k,r-k+1)$ density, with
mode $(k-1)/(r-1)$ when $k>1$ and total mass $wm(1-k/r)$. The
release occurs near the code threshold, and fresh questions continue
to add entropy at marginal leverage one afterwards, which creates the
declining side of the peak. At $k = r$ the code is the full product
space, $\tau_{r,r}\equiv1$, and the deduction density vanishes
identically.

\subsection{A lapsed MDS block}

Let $\mathcal K + y\subset\mathbb F_{A}^r$ be an MDS coset with fixed
offset $y$. For $0\le\delta\le1$ draw uniformly from the coset with
probability $1-\delta$ and uniformly from all $A^r$ answer strings
with probability $\delta$, which has full support for every
$0<\delta\le1$. For $j\le k$ both mixture components project to the
uniform distribution on $\mathbb F_{A}^j$, so $H_\delta(j) = jm$. For
$j\ge k$ the projection of the code has $A^k$ consistent patterns,
each of probability
$(1-\delta)A^{-k}+\delta A^{-j} = A^{-k}\big[(1-\delta)+\delta
A^{k-j}\big]$, while the other $A^j-A^k$ patterns each have
probability $\delta A^{-j}$. Summing the two entropy contributions,
%
\begin{equation}
  H_\delta(j)
  = \Lambda\big[km-\log_2\Lambda\big]
  + \delta\big(1-A^{\,k-j}\big)\big[jm-\log_2\delta\big],
  \qquad
  \Lambda := (1-\delta)+\delta A^{\,k-j},
  \label{eq:applapsed}
\end{equation}
%
for $j\ge k$, with $\eta_\delta(i) = H_\delta(i+1)-H_\delta(i)$. It
equals $m$ below threshold; above threshold it declines toward a
positive lapse tail, and the exact finite differences should be used
rather than replacing the whole tail immediately by $\delta m$.

\section{Local noise as a term-by-term
perturbation}\label{app:noise}

\subsection{A general channel construction}

Let $P_{\mathcal V}$ be any law on
$\mathcal A^{|\mathcal V|}$ and, for $0\le\delta\le1$, let
$\Pr_\delta(y\mid a)$ be a full-support channel on the answer
alphabet tending to the identity as $\delta\to0$. Applying it
independently to each coordinate,
%
\begin{equation}
  P_{\mathcal V,\delta}(y_{\mathcal V})
  = \sum_{a_{\mathcal V}\in\mathcal A^{|\mathcal V|}}
  P_{\mathcal V}(a_{\mathcal V})
  \prod_{q\in\mathcal V}\Pr\nolimits_\delta(y_q\mid a_q),
  \label{eq:appchannel}
\end{equation}
%
or alternatively using the lapse \eqref{eq:lapse}, the full map law
remains a product over blocks and has full support if each noisy
block law does.

\subsection{Why vanishing local noise preserves the limit}

For fixed $|\mathcal V|$ the block probability simplex is
finite-dimensional, so
$\|P_{\mathcal V,\delta}-P_{\mathcal V}\|_{\rm TV}\to0$ as
$\delta\to0$. Entropy is continuous on a finite alphabet, so every
subset entropy and increment converges,
$H_{\mathcal V,\delta}(j)\to H_{\mathcal V}(j)$ and
$\eta_{\mathcal V,\delta}(i)\to\eta_{\mathcal V}(i)$. The finite
hypergeometric weights, and their limiting Bernstein weights, are
nonnegative and sum to one, so writing $\gamma_{{\rm noisy},n}$ and
$\gamma_{{\rm hard},n}$ for the profiles computed from the noisy and
hard coefficients at size $n$,
%
\begin{equation}
  \sup_{t\in[0,1]}
  \big|\gamma_{{\rm noisy},n}(t)-\gamma_{{\rm hard},n}(t)\big|
  \;\le\;
  \sum_{\mathcal V\in\mathcal P_n}\frac{|\mathcal V|}{Q}
  \max_i\big|\eta_{\mathcal V,\delta_{\mathcal V,n}}(i)
  -\eta_{\mathcal V,0}(i)\big|
  \;\longrightarrow\;0 .
  \label{eq:appnoisebound}
\end{equation}
%
The integrals converge uniformly as well, and wherever the limiting
received density is bounded away from zero, leverage converges. For
fixed block sizes continuity makes the last limit automatic if
$\max_{\mathcal V}\delta_{\mathcal V,n}\to0$; the aggregate condition
is the more general statement, and for varying block sizes a uniform
condition is needed, model-specific ones such as $r(n)\delta_n\to0$
for parity blocks being sharper. At fixed $\delta>0$ none of these
perturbations disappears with $n$: the exact noisy profile remains the
thermodynamic limit.

\section{Arbitrary mixtures and the cocktail}\label{app:mixtures}

\subsection{Symbolic mixture law}

Partition the questions into independent component families, type
$\kappa$ occupying fraction $w_\kappa$ with initial one-question
entropy $h_{0,\kappa}$, macroscopic profile $\gamma_\kappa$ and
accumulated received density $g_\kappa = \int_0^t\gamma_\kappa$.
Entropy adds across independent components, so $h_0$, $\gamma$ and
$g$ are the corresponding $w_\kappa$-weighted sums. Since
$h_{0,\kappa}-(1-t)\gamma_\kappa(t) = L_\kappa(t)g_\kappa(t)$ for
each component, summing gives \eqref{eq:mixture}: the mixture weights
in leverage are dynamic received-entropy shares
$w_\kappa g_\kappa(t)$, not the fixed question shares $w_\kappa$.

\subsection{Coin component of the cocktail}

For the two-coin component take $(\theta_1,\theta_2) = (0.2,0.7)$
with weights $(\tfrac14,\tfrac34)$ and let the same scalar latent
bias govern all bits. With $P_{s,z}$ and $F_s$ as in
\eqref{eq:coinPlk} and \eqref{eq:Fs}, the initial entropy per bit for
an $m = 4$-bit answer is
%
\begin{equation}
  \frac{F_4}{4} = \frac{3.793149840}{4} = 0.948287460,
  \label{eq:appcoinF4}
\end{equation}
%
which is \emph{not} $h_2(0.575) = 0.983708263$, because marginalizing
the shared latent makes the four bits dependent. After a
vanishing-fraction identification layer the entropy per bit is
%
\begin{equation}
  h_{\rm cond,coin}
  = \tfrac14 h_2(0.2) + \tfrac34 h_2(0.7) = 0.841450198 .
  \label{eq:appcoincond}
\end{equation}

\subsection{Lapsed-code and fresh components}

For $A = 16$, $r = 16$, $k = 4$ and $\delta = 0.1$, compute
$H_\delta(j)$ from \eqref{eq:applapsed}, set
$\eta_\delta(i) = H_\delta(i+1)-H_\delta(i)$, and normalize the code
profile per bit,
%
\begin{equation}
  \gamma_{\rm lapse}(t)
  = \frac14\sum_{i=0}^{15}\beta_{15,i}(t)\,\eta_\delta(i),
  \qquad
  \int_0^1\gamma_{\rm lapse}
  = \frac{H_\delta(16)}{16\cdot4} = 0.332328056 .
  \label{eq:applapseprofile}
\end{equation}
%
The fresh component consists of independent bits with bias $0.3$, so
its per-bit profile and initial entropy are both $h_2(0.3)$.

\subsection{Assembling the cocktail}

With question fractions $\tfrac14,\tfrac12,\tfrac14$ the initial
per-bit entropy is \eqref{eq:cocktailh0}, the profile for $t>0$ is
\eqref{eq:cocktailgamma}, and
%
\begin{equation}
  g_{\rm cocktail}(t)
  = \tfrac14 h_{\rm cond,coin}t
  + \tfrac12\int_0^t\gamma_{\rm lapse}
  + \tfrac14 h_2(0.3)\,t,
  \label{eq:appcocktailg}
\end{equation}
%
the scale factor $m$ cancelling from
$L = [h_{\rm cocktail}-(1-t)\gamma_{\rm cocktail}]/g_{\rm cocktail}$.
At the origin the profile's right limit is
$\gamma_{\rm cocktail}(0^+) = \tfrac14(0.841450198)+\tfrac12
+\tfrac14h_2(0.3) = 0.930685274$, so \eqref{eq:boundarylayer} gives
the exact singular coefficient
%
\begin{equation}
  \lim_{t\downarrow0}tL(t)
  = \frac{h_{\rm cocktail}-\gamma_{\rm cocktail}(0^+)}
         {\gamma_{\rm cocktail}(0^+)}
  \simeq 0.02870 ,
  \label{eq:appcocktailsing}
\end{equation}
%
while at the endpoint
$g_{\rm cocktail}(1) = \tfrac14(0.841450198)
+\tfrac12(0.332328056)+\tfrac14h_2(0.3) = 0.596849302$ and
$L(1) = 0.957394590/0.596849302 \simeq 1.60408$. Numerical
evaluation of the same closed form gives the dip and peak quoted in
Section~\ref{sec:zoo}.

\section{What can be inferred from an observed
curve?}\label{app:inverse}

If $h_0$ and an ideal macroscopic leverage curve $L(t)$ are known,
then $g'(t) = \gamma(t)$ and \eqref{eq:master} imply the linear
differential equation
%
\begin{equation}
  (1-t)g'(t) + L(t)g(t) = h_0,
  \qquad g(0) = 0,
  \label{eq:appinverse}
\end{equation}
%
and for an admissible $L$, solving it for the regular, finite,
nonnegative solution recovers the entropy profile $\gamma = g'$. The
qualification matters when $L$ has a $1/t$ singularity, since
$g(0) = 0$ is then a singular initial condition rather than a
standard nonsingular initial-value problem.

The inverse map does not recover a unique prior. Many microscopic
distributions share a subset-entropy profile, and finite data
introduce further ambiguity; for a trained model, question selection,
optimization dynamics, misspecification and calibration also affect
the measured curve. The defensible inference is therefore about broad
timing classes: a $1/t$ term indicates information learned on a
subextensive clock; a flat curve with $L>1$ is consistent with
first-touch duplication, while $L\equiv1$ is also the independent
baseline; a rising curve is consistent with late constraints; and an
interior peak is consistent with threshold release followed by
continued reception. This is useful structural evidence, not an
identification theorem for the real-world prior.

\section{Endpoints, jumps, and boundary
layers}\label{app:endpoints}

At the end of the run,
$L(1) = h_0/g(1) = \lim_{n\to\infty} Q G_{n,1}/G_{n,Q}$. At the
origin there are two different cases. If $h_0 = \gamma(0^+)>0$ and
$\gamma$ is right-differentiable, first-order expansion gives
%
\begin{equation}
  L(0^+) = 1-\frac{\gamma'(0^+)}{\gamma(0^+)} .
  \label{eq:appLzero}
\end{equation}
%
If instead $h_0>\gamma(0^+)>0$, an initial boundary layer has
disappeared into $t = 0$ and
%
\begin{equation}
  L(t)\;\sim\;\frac{h_0-\gamma(0^+)}{\gamma(0^+)}\,\frac1t
  \qquad (t\downarrow0), 
  \label{eq:boundarylayer}
\end{equation}
%
which is the source of the spike and coin-mixture hyperbolas.

A separate caveat concerns interior thresholds. If $\gamma$ merely
reaches zero continuously, no instantaneous avalanche follows; a
macroscopic jump occurs when $\gamma(t_*^-)>\gamma(t_*^+)$, and
especially when the profile jumps to zero. The remaining-entropy
density then drops by
$(1-t_*)\big[\gamma(t_*^-)-\gamma(t_*^+)\big]$ and cumulative
leverage can jump as well; the local per-step leverage inside the
shrinking transition layer may diverge while the cumulative $L(t)$
need not. Finally, if $G_{n,Q} = o(Q)$ the normalization of
Section~\ref{sec:tdlcalc} erases the entire received channel, and one
should then introduce a microscopic or mesoscopic time scale rather
than conclude automatically that leverage has a meaningful infinite
value.